\section{Agent Swarms：多 Agent 并行协作}

\subsection{Agent Swarm 范式}

单 agent 模式面临任务复杂度的瓶颈。Kimi 提出了 \textbf{Agent Swarm} 范式：

\begin{itemize}
    \item 一个\textbf{主 agent（orchestrator）}负责任务编排
    \item 主 agent 可以生成一组\textbf{子 agent}，分配子任务
    \item 子 agent 并行执行，主 agent 收集结果
    \item 整个过程可迭代进行
\end{itemize}

\begin{knowledgebox}{类比人类组织}
Agent Swarm 类似于一家公司的组织结构：CEO 负责分解和分配任务（orchestrator），AI 研究员、Web 开发者、物理研究员等各司其职（sub-agents），最终由 fact checker 等角色汇总结果。
\end{knowledgebox}

实验表明，Agent Swarm 能显著降低任务执行时间，尤其在高复杂度任务上效果突出。当扩展到 100 甚至 1000 个子 agent 时，可以在可接受的时间内完成复杂任务。

\subsection{Agent Swarm 的 RL 训练}

\begin{importantbox}{三项奖励函数}
除了标准的结果奖励（outcome reward）外，Agent Swarm 的 RL 训练引入了两个额外的奖励信号：
\begin{enumerate}
    \item \textbf{Instantiation Reward}：激励子 agent 的并行实例化，防止"串行坍缩"——即模型退化为单 agent 执行
    \item \textbf{Finish Reward}：确保子任务有较高的完成率，防止模型"刷"第一个奖励——只生成大量子 agent 但不完成任务
    \item \textbf{Outcome Reward}：标准的任务完成度奖励
\end{enumerate}
前两个奖励在训练过程中逐渐衰减权重（decay strategy），早期鼓励探索，后期侧重结果。
\end{importantbox}

\begin{warningbox}{Agent Swarm 的训练陷阱}
如果只有 outcome reward，模型容易退化为单 agent 模式（串行坍缩）。如果只加 instantiation reward，模型可能 hack 这一信号——生成大量伪子任务但从不完成。必须同时使用三种奖励信号才能实现稳定训练。
\end{warningbox}

\subsection{本章小结}

Agent Swarm 通过多 agent 并行协作开辟了新的 scaling 维度。三项精心设计的奖励函数配合衰减策略，使模型能有效学习并行编排和任务分解能力。


